\documentclass[12pt,a4paper]{article}

\usepackage[utf8]{inputenc}
\usepackage[T1]{fontenc}
\usepackage[polish]{babel}
\usepackage{amsmath}
\usepackage{amssymb}
\usepackage{geometry}
\usepackage{array}
\usepackage{listings}
\usepackage{needspace}

\geometry{margin=2.5cm}

% Naglowki sekcji nie powinny zostawac same na koncu strony.
\newcommand{\sekcja}[1]{%
  \Needspace{8\baselineskip}%
  \section*{#1}%
  \nopagebreak[4]%
}

% Blok naglowek + krotki wstep trzyma razem na jednej stronie.
\newcommand{\sekcjazwstep}[2]{%
  \Needspace{10\baselineskip}%
  \section*{#1}%
  \nopagebreak[4]%
  #2%
}

\widowpenalty=10000
\clubpenalty=10000

\lstset{
    basicstyle=\ttfamily\small,
    breaklines=true,
    frame=single,
    columns=fullflexible,
    keepspaces=true,
    showstringspaces=false,
    literate={ą}{{\k{a}}}1 {ć}{{\'c}}1 {ę}{{\k{e}}}1 {ł}{{\l}}1 {ń}{{\'n}}1 {ó}{{\'o}}1 {ś}{{\'s}}1 {ź}{{\'z}}1 {ż}{{\.z}}1 {Ą}{{\k{A}}}1 {Ć}{{\'C}}1 {Ę}{{\k{E}}}1 {Ł}{{\L}}1 {Ń}{{\'N}}1 {Ó}{{\'O}}1 {Ś}{{\'S}}1 {Ź}{{\'Z}}1 {Ż}{{\.Z}}1 % Obsługa polskich znaków w kodzie
}

\title{Podstawy Symulacji Komputerowej\\Lista 0 -- zadanie 2}
\author{Ryszard Rębowski}
\date{}

\begin{document}

\maketitle

\sekcjazwstep{Treść zadania}{%
Wiadomo, że dla zmiennej losowej $X$ mamy:
\[ X(\Omega)=\{1,2,3,4\} \]
\[ P(\{\omega\in\Omega:X(\omega)=i\})=ic, \quad \text{dla pewnej } c\in \mathbb{R}. \]
Obliczyć dwoma metodami $P(\{\omega\in\Omega:2\le X(\omega)\le3\})$.
}

\sekcjazwstep{Rozwiązanie analityczne}{%
Z własności rozkładu prawdopodobieństwa (warunek unormowania) wiemy, że suma prawdopodobieństw wszystkich zdarzeń elementarnych musi wynosić $1$. Zatem:
\[ \sum_{i=1}^{4} P(X=i) = 1 \]
\[ 1c + 2c + 3c + 4c = 1 \]
\[ 10c = 1 \implies c = \frac{1}{10}. \]

Mając stałą $c$, rozkład zmiennej losowej $X$ prezentuje się następująco:
$P(X=1)=0{,}1$, $P(X=2)=0{,}2$, $P(X=3)=0{,}3$, $P(X=4)=0{,}4$.

\textbf{Metoda 1: Bezpośrednie sumowanie prawdopodobieństw} \\
Z racji, że interesują nas zdarzenia dla $X=2$ oraz $X=3$, dodajemy bezpośrednio ich prawdopodobieństwa z wyznaczonego rozkładu:
\[ P(2 \le X \le 3) = P(X=2) + P(X=3) = 0{,}2 + 0{,}3 = 0{,}5. \]

\textbf{Metoda 2: Sumowanie z definicji podanej w zadaniu} \\
Wykorzystujemy podany wzór $P(X=i) = i \cdot c$:
\[ P(2 \le X \le 3) = \sum_{i=2}^{3} i \cdot c = 2c + 3c = 5c. \]
Podstawiając $c = \frac{1}{10}$, otrzymujemy:
\[ 5 \cdot \frac{1}{10} = \frac{5}{10} = 0{,}5. \]
}

\sekcjazwstep{Implementacja}{%
Poniżej przedstawiono kod w języku Python symulujący rozwiązanie tego zadania wraz z generacją wykresów rozkładu zmiennej losowej i jej dystrybuanty.
}

\Needspace{22\baselineskip}
\begin{lstlisting}[language=Python]
import matplotlib.pyplot as plt
import numpy as np

print("\n1.Znalezienie stalej c:")
print("X(Ω) = {1, 2, 3, 4}")
print("P(X = i) = i·c")
print("\nWarunek: suma wszystkich prawdopodobienstw = 1")
print("P(X=1) + P(X=2) + P(X=3) + P(X=4) = 1")
print("1·c + 2·c + 3·c + 4·c = 1")
print("10·c = 1\n")

c = 1/10
print(f"\033[1mStala c:\033[0m {c} \n")
values = [1, 2, 3, 4]
probabilities = [i * c for i in values]

print(f"P(X=1) = 1·{c} = {probabilities[0]}")
print(f"P(X=2) = 2·{c} = {probabilities[1]}")
print(f"P(X=3) = 3·{c} = {probabilities[2]}")
print(f"P(X=4) = 4·{c} = {probabilities[3]} \n")
print(f"\033[1mSuma:\033[0m {sum(probabilities)} \n")

print("Metoda 1: Bezposrednie sumowanie")
print("P(2 ≤ X ≤ 3) = P(X=2) + P(X=3)")
print(f"= 2·c + 3·c")
print(f"= 2·{c} + 3·{c}")
print(f"= {2*c} + {3*c}")

result_method1 = probabilities[1] + probabilities[2]
print(f"\033[1m= {result_method1} \033[0m\n")

print("Metoda 2: Przez sumowanie zdan 2 i 3 z definicji")
print("P(2 ≤ X ≤ 3) = Σ P(X=i) dla i ∈ {2, 3}")

result_method2 = sum(probabilities[i-1] for i in [2, 3])
print(f"= {probabilities[1]} + {probabilities[2]}")

print(f"\033[1m= {result_method2} \033[0m\n")

print("Weryfikacaja zgodnosci wynikow:")
print(f"\033[1mMetoda 1: {result_method1}\033[0m")
print(f"\033[1mMetoda 2: {result_method2}\033[0m")
print(f"\033[1mZgodnosc: {'Tak' if abs(result_method1 - result_method2) < 1e-10 else 'Nie'}\033[0m\n")

fig, axes = plt.subplots(1, 2, figsize=(14, 5))

ax1 = axes[0]
colors = ['lightblue', 'salmon', 'salmon', 'lightblue']
bars = ax1.bar(values, probabilities, color=colors, edgecolor='black', linewidth=2, width=0.6)

bars[1].set_color('salmon')
bars[2].set_color('salmon')

ax1.set_xlabel('X', fontsize=12, fontweight='bold')
ax1.set_ylabel('P(X = i)', fontsize=12, fontweight='bold')
ax1.set_title('Rozklad zmiennej losowej X', fontsize=13, fontweight='bold')
ax1.set_xticks(values)
ax1.grid(axis='y', alpha=0.3, linestyle='--')
ax1.set_ylim([0, max(probabilities) * 1.2])

for i, (val, prob) in enumerate(zip(values, probabilities)):
    ax1.text(val, prob + 0.01, f'{prob:.1f}', ha='center', fontsize=11, fontweight='bold')

ax1.text(0.5, 0.95, 'Kolor jasnoniebieski: nie wliczone do P(2 ≤ X ≤ 3)', 
         transform=ax1.transAxes, fontsize=10, verticalalignment='top',
         bbox=dict(boxstyle='round', facecolor='lightblue', alpha=0.5))
ax1.text(0.5, 0.85, 'Kolor lososiowy: wliczone do P(2 ≤ X ≤ 3)', 
         transform=ax1.transAxes, fontsize=10, verticalalignment='top',
         bbox=dict(boxstyle='round', facecolor='salmon', alpha=0.5))

ax2 = axes[1]
cumulative_probs = np.cumsum(probabilities)
x_cdf = [0] + values
y_cdf = [0] + list(cumulative_probs)

ax2.step(x_cdf, y_cdf, 'b-', linewidth=2, marker='o', markersize=8, where='post', label='F(x) = P(X ≤ x)')
ax2.fill_between([1.99, 3.01], 0, 1, alpha=0.3, color='salmon', label='P(2 ≤ X ≤ 3)')

ax2.set_xlabel('X', fontsize=12, fontweight='bold')
ax2.set_ylabel('F(X)', fontsize=12, fontweight='bold')
ax2.set_title('Dystrybuanta zmiennej losowej X', fontsize=13, fontweight='bold')
ax2.grid(True, alpha=0.3, linestyle='--')
ax2.set_xlim([-0.5, 5])
ax2.set_ylim([-0.05, 1.1])
ax2.legend(fontsize=10)

for i, cum_prob in enumerate(cumulative_probs):
    ax2.text(x_cdf[i], cum_prob + 0.05, f'{cum_prob:.2f}', ha='center', fontsize=9)

plt.tight_layout()
print("\nWizualizacja zapisana do: zad2_lista0_visualization.png \n")
plt.show()

print(f"\033[1mPodsumowanie:\033[0m")
print(f"P(2 ≤ X ≤ 3) = {result_method1} = {result_method1.as_integer_ratio()[0]}/{result_method1.as_integer_ratio()[1]} = 50%")
\end{lstlisting}

\sekcjazwstep{Odpowiedź końcowa}{%
Obydwie metody analityczne, jak i weryfikacja numeryczna za pomocą skryptu Python, doprowadziły do tego samego wyniku. Szukane prawdopodobieństwo wynosi:
}

\Needspace{12\baselineskip}
\[
\boxed{P(2 \le X \le 3) = 0{,}5}
\]

czyli:

\[
\boxed{50\%}.
\]

\end{document}